%% Rewritten 2026-09-08 after the Marvin 26.1 regeneration and the reversibility
%% overhaul. Every number here is measured against the shipped files; see the
%% inline comments for the command that reproduces each. Detail in
%% Supplementary Methods S2.
\subsection{Multi-source thermodynamics}\label{sec:methods-thermo}

The 2020 release combined two sources: the historical group contribution approach~\cite{jankowski2008} and eQuilibrator~\cite{flamholz2012}. This update carries four sources, each stored as its own record: group contribution, eQuilibrator~3.0~\cite{beber2022}, dGPredictor~\cite{wang2021}, and experimental values where they exist. Only the experimental values carry a published direction. The three predictors were released as energy estimators, so every direction reported here is an inference from their energies. All values are reported at pH~7.0, ionic strength 0.25\,M, pMg~3.0 and 298.15\,K; the basis for each condition and its effect on the released energies are given in Supplementary Methods~S2.

\paragraph{Protonation.} eQuilibrator's Legendre transform requires a \emph{macroscopic} pKa ladder per compound: an ordered list in which consecutive species differ by one proton. The count of entries around the desired pH sets the proton count of the major microspecies, so the composition of the ladder changes the transformed energy directly. We generate the pKa ladders using ChemAxon Marvin~\cite{marvin} for use with eQuilibrator, but we're conscious that these data cannot be regenerated without a license, so we also include the pKas generated by an open-source alternative, MolGpKa~\cite{pan2021}, for every compound with a structure, and the measured ladders of Alberty~\cite{alberty2003}, versioned per structure source in an extended pKa dictionary.

\paragraph{Experimental anchors.} eQuilibrator and dGPredictor both draw on the NIST Thermodynamics of Enzyme-catalyzed Reactions (TECR) Database~\cite{goldberg2004}, which is keyed to KEGG and reaches other namespaces through MetaNetX mappings that conflict structurally with ModelSEED identifiers. We therefore rebuilt the anchors, curating each conflict so that every TECR value maps directly onto a ModelSEED reaction, and evaluated openTECR (\url{https://github.com/opentecr}), a communal re-curation of the same source, to extend the set.

\paragraph{Uncertainty.} The three sources each adopt a different approach for estimating the uncertainty surrounding the predicted energy of reaction, so the scale of uncertainty is not directly comparable between them (Figure~\ref{fig:thermo}C). We therefore calibrate each source's reported uncertainty against the experimental anchors rather than reading it at face value; the resulting medians and the direction of each source's bias are reported with the results.

\paragraph{Direction.} For historical reasons, we use the same heuristics for deriving reaction direction from the Group contribution approach~\cite{jankowski2008} unchanged from 2020, so that series stays comparable across releases. For eQuilibrator and dGPredictor we use one rule set built on the reversibility index from Noor \textit{et al.}~\cite{noor2012}: a reaction is called irreversible when $|\ln\Gamma|$ exceeds $\ln(1000)$ by more than one propagated standard deviation~\cite{beber2022}, $\Gamma$ being the fold-change of reactant concentration required to change reaction direction). However, we extend the work by Noor \textit{et al.} to require the uncertainty to clear the threshold as well. We report fewer directions than the published index would on the same energies, and we separate \emph{reversible} reactions from \emph{undetermined} cases. An exception we implement is for reactions whose uncertainty is small relative to the threshold and yet still straddles it: meaning the estimate lies within one standard deviation of the threshold, and these we report as reversible.

\paragraph{Transport.} Compartments are \emph{relative} indices --- 0 inside, 1 outside --- not named compartments, so one reaction can be instantiated against a bacterial membrane, a mitochondrion or a plastid envelope when a template builds a model. Because the database is not committed to a cell type, transport reactions are scored from stoichiometry and energy alone: neither membrane potential nor pH gradient enters the estimate. The consequence is specific. A primary active transporter carries ATP hydrolysis in its stoichiometry, every predictor estimates that tightly, and the reaction earns a confident grade from chemistry that was never in question --- 96\% of gold-graded transport reactions are ATP-coupled and only 2.7\% translocate a proton. Such grades measure confidence in the coupled chemistry, not evidence that the transporter runs that way in a given organism; every reaction carries an \texttt{is\_transport} field so that this can be accounted for downstream. Reactions whose only change is of compartment carry no net chemistry and none is graded gold. dGPredictor commits on 1.4\% of transport reactions against 17.3\% elsewhere, so transport grades also rest on fewer independent sources.
